\documentclass[11pt,a4paper]{article}

\usepackage[T1]{fontenc}
\usepackage[utf8]{inputenc}
\usepackage[margin=1in]{geometry}
\usepackage{amsmath}
\usepackage{amssymb}
\usepackage{booktabs}
\usepackage{microtype}
\usepackage[colorlinks=true,linkcolor=blue,urlcolor=blue]{hyperref}

\newcommand{\epsm}{\varepsilon_{\mathrm{machine}}}

\title{Eigenvalue Solvers versus Linear Solvers\\
       \large A Comparison Based on Trefethen \& Bau,
       \emph{Numerical Linear Algebra} (SIAM, 1997)}
\author{}
\date{}

\begin{document}

\maketitle

\noindent
All page numbers refer to the printed pagination of the book. In the scanned
file \path|docs/Bau_Numerical_linear_Algebra_BOOK.pdf|, the PDF page
number is the printed page number plus~10.

\bigskip

\noindent
Here is what the book actually says when it puts the two side by side. The
comparison runs along five axes.

\section{The fundamental difference: finite vs.\ iterative}

This is the book's headline point, made twice --- Lecture~25 (p.~192) and again
in the Appendix (pp.~323--324).

The system $Ax = b$ is solvable by a \emph{finite} algorithm. Gaussian
elimination would produce the exact answer in a fixed, known number of steps if
arithmetic were exact. Eigenvalues cannot be: any polynomial rootfinding problem
\emph{is} an eigenvalue problem (via the companion matrix, eq.~25.3), and
Abel/Galois says degree $\geq 5$ roots have no closed form. Hence:

\begin{quote}
``Any eigenvalue solver must be iterative.'' (p.~192)
\end{quote}

The Appendix makes this the dividing line of the whole field: Gaussian
elimination is the one premier example where ``operation counts, machine
architectures, and rounding errors --- that's all you have to understand.''
Eigenvalue problems, like everything else with a nonlinearity in it, need a
theory of \emph{convergence} on top.

Trefethen's framing (Figure~32.1, p.~247): a direct linear solver ``makes no
progress at all until $O(m^3)$ operations are completed''; an iterative method
converges geometrically toward $\epsm$.

\section{Cost: the gap is a small constant, not an order}

Despite being ``unsolvable in principle,'' eigenvalues are barely more
expensive. From p.~193: in practice the eigenvalue computation ``differs from
the solution of linear systems by only a small constant factor, typically closer
to 1 than 10.'' The $\varepsilon$-dependence of the iteration is only
$\log\!\left(\lvert \log \epsm \rvert\right)$ (Exercise~25.2).

Leading-term flop counts, all from the book:

\begin{center}
\begin{tabular}{@{}llc@{}}
\toprule
\textbf{Task} & \textbf{Work} & \textbf{Reference} \\
\midrule
Cholesky factorization (SPD solve)   & $\sim m^3/3$    & (23.4), p.~176 \\
Gaussian elimination / LU solve      & $\sim 2m^3/3$   & (20.4), p.~152 \\
Householder QR (square)              & $\sim 4m^3/3$   & (10.9), p.~75  \\
Tridiagonal reduction (hermitian)    & $\sim 4m^3/3$   & (26.2), p.~200 \\
Hessenberg reduction (general)       & $\sim 10m^3/3$  & (26.1), p.~199 \\
Symmetric eigenvalues only           & $\sim 4m^3/3$   & p.~223 \\
Symmetric eigenvalues $+$ vectors, QR algorithm
                                     & $\sim 9m^3$     & p.~232 \\
Symmetric eigenvalues $+$ vectors, divide \& conquer
                                     & $\sim 4m^3$, often $\sim 8m^3/3$ & p.~232 \\
Matrix--matrix product (for scale)   & $2m^3$          & p.~223 \\
\bottomrule
\end{tabular}
\end{center}

So the direct comparison the book invites: symmetric eigenvalues cost
$4\times$ a Cholesky, $2\times$ an LU, and exactly the same as a Householder
QR --- and are ``two-thirds the cost of computing the product of a pair of
$m \times m$ matrices'' (p.~223), which the authors call ``not a bad result at
all.''

The cost of \emph{eigenvectors} is where the two genuinely diverge: $9m^3$ vs
$2m^3/3$ is more than an order of magnitude.

\section{The structural inversion: which phase dominates}

This is the most interesting asymmetry, p.~194. For a linear solve, all the work
is in the finite factorization; the triangular back-substitutions are $O(m^2)$,
negligible.

For a symmetric eigenproblem it is the same shape but for the opposite reason:
Phase~1 (reduction to tridiagonal) is finite and $O(m^3)$; Phase~2 (the formally
infinite QR iteration) is $O(m)$ iterations $\times$ $O(m)$ flops $= O(m^2)$.
The book calls this paradoxical:

\begin{quote}
``the `infinite' part of the algorithm is in practice not merely as fast as the
`finite' part, but an order of magnitude faster.''
\end{quote}

Practical consequence: a symmetric eigensolve is dominated by a finite, direct,
factorization-like step, exactly as a linear solve is. The iteration is
asymptotically free. Note also that Phase~1 is not optional --- skipping it
makes each Phase-2 step $O(m^3)$ and the total $O(m^4)$.

\section{Conditioning: the eigenvalue problem is better behaved (if symmetric)}

Here the comparison reverses in the eigensolver's favor.

\begin{itemize}
  \item \textbf{Linear system.} The condition number is
        $\kappa(A) = \lVert A \rVert \, \lVert A^{-1} \rVert$
        (Theorem~12.2, p.~95). You ``must always expect to lose
        $\log_{10} \kappa(A)$ digits.'' Ill-conditioning is unbounded and
        routine.

  \item \textbf{Symmetric/normal eigenvalues.}
        $\lvert \tilde{\lambda}_j - \lambda_k \rvert \leq \lVert \delta A \rVert_2$
        (eq.~26.4, Exercise~26.3b, p.~201). The absolute condition number is
        $1$. Combined with backward stability of tridiagonal reduction
        (Thm.~26.1) and QR (Thm.~29.1), computed eigenvalues satisfy
        $\lvert \tilde{\lambda}_j - \lambda_j \rvert
         = O(\epsm) \cdot \lVert A \rVert$ --- no $\kappa(A)$ anywhere.

  \item \textbf{Nonsymmetric eigenvalues.} The advantage evaporates.
        Bauer--Fike gives disks of radius
        $\kappa(V) \lVert \delta A \rVert_2$ --- the condition number of the
        \emph{eigenvector} matrix enters, and the book spends Exercise~26.1 on
        pseudospectra because that bound can be badly pessimistic or badly
        violated in effect.
\end{itemize}

\section{Large and sparse: the two converge}

In Part~VI the distinction largely dissolves (pp.~244--248). Both problem
classes are attacked by Krylov methods through the same $x \mapsto Ax$ black
box, at $O(\nu m)$ per matrix-vector product instead of $O(m^2)$. Both are
governed by the same cost decomposition --- \emph{number of steps} $\times$
\emph{work per step} --- where dense direct algorithms sit at $O(m)$ steps
$\times$ $O(m^2)$, and ``the ideal iterative method reduces the number of steps
from $m$ to $O(1)$ and the work per step from $O(m^2)$ to $O(m)$.''

The paired algorithms are explicitly analogous (p.~246): CG solves an SPD system
fast when eigenvalues are clustered away from the origin; Lanczos finds
eigenvalues fast when those eigenvalues are well separated from the rest of the
spectrum. In both cases convergence is a question of approximation theory ---
polynomial approximation on the spectrum --- not of the problem type.

\bigskip

\noindent
\textbf{Caveat.} None of this covers shift-invert generalized eigensolves such
as \path|eigsh(K, M, sigma=...)|, where each Lanczos step contains a sparse
solve against $K - \sigma M$. There the eigensolver's cost is literally a
multiple of a linear solver's, and the book's clean separation does not apply.
That is \path|docs/ARPACK.pdf| territory.

\end{document}
